\subsection{The computational object beyond matrix diagnostics}

The method addresses a question that matrix amplitude and conditioning leave
open. The quantities $\mu_J$, $\rho_1$, $\rho_g$, and $\kappa(J)$ quantify
representation-dependent change or worst-direction sensitivity. They do not
identify whether the selected solution activates the perturbation range.
Proposition~\ref{prop:p2} makes that information deficit exact: all four
matrix-level quantities can be matched while the realized response changes
from the guarded-inverse bound to zero. Structured and componentwise
conditioning sharpen the admissible perturbation model
\cite{S39,S40,S41}, but a selected response still requires the direction
$Ez$ and its propagation by $A^{-1}$.

The second advance concerns computation of that response. Direct subtraction
and the response equation are algebraically equivalent in exact arithmetic,
yet they form different finite-precision operation graphs.
Propositions~\ref{prop:p4} and \ref{prop:p5} provide separate a posteriori
error enclosures for those graphs. This is the central difference from a
generic verified solve: the solve-error ingredients are propagated to a
specific perturbation-response object, and the forcing-construction defect is
retained when the declared perturbation and realized solve operands differ at
the last bit. The threshold relation in Corollary~\ref{cor:threshold} then
turns the enclosed response into an above, below, or unresolved result without
using a condition number as a surrogate decision.

\subsection{Why the enclosures are conservative}

Complete coverage and useful sharpness are different properties. The
effectivity factorization in Eq.~\eqref{eq:effectivity_factorization}
localizes the dominant source of conservativeness. Median inflation caused by
the Frobenius majorant was almost unity, whereas replacing the realized
residual direction by the worst inverse direction contributed median factors
between $1.36\times10^3$ and $2.31\times10^4$. The large central ratios are
therefore not primarily a Frobenius-versus-spectral artifact; they arise
because a valid normwise residual bound discards directional alignment.

The multivariable analysis also prevents overinterpreting a single marginal
correlation. Condition exponent explained most registered variation in the
development design, but its drop-one partial $R^2$ fell from $0.855$ to
$0.029$ on the original holdout. Matrix family, perturbation amplitude, and
reference lane likewise changed importance across the two finite grids. The
data support a mechanism---worst-direction replacement can be highly
conservative---but not a universal effectivity scaling law. The extreme
$\log_{10}I_{\rm eff}=86.90$ further illustrates why ratios become
unstable when the true error approaches zero even though absolute coverage
remains valid.

\subsection{What the two routes reveal}

The response equation reduced decimal-source reference error in 295/300
development cases and 271/300 new-holdout cases. Its enclosure was also
smaller in 283/300 and 229/300 cases. Binary64-operand comparisons were less
separated: the response route had lower error in 199/300 development and
160/300 holdout cases, and lower bounds in 229/300 and 188/300. These
differences are consistent with the operation graphs. Direct subtraction can
lose small responses between two complete parent solutions; the response
equation avoids that subtraction but inherits forcing-construction and
response-solve errors. The new result is not that one graph always wins, but
that both can now be compared through valid route-specific enclosures rather
than raw arithmetic error alone.

\subsection{Computational-mechanics meaning and scope}

For a frozen coupled tangent, the method answers whether the correction
induced by a tangent-only modification is above, below, or unresolved with
respect to an independently supplied response threshold. A solver developer
may connect that threshold to an accepted Newton-correction tolerance, a
residual scale, an energy norm, or a quantity-of-interest budget. Such a
choice converts a vague statement that the matrix change is \emph{small}
into a testable statement about the selected correction at one nonlinear
state. M8 and M16 demonstrate this chain for
displacement--pressure--saturation operators with explicit field scales; both
operation graphs excluded the special threshold zero.

The present implementation is intentionally a high-precision reference
benchmark. It does not select a perturbation, prove nonlinear convergence,
or certify an engineering model. It is also not an outward-rounded interval
solver. Sparse estimation, controlled large-system sequences, blind transfer,
and coverage--tightness--cost validation belong to a separate scalable-method
study. Consolidating these limits here preserves the positive result: a
coordinate-consistent, operation-graph-specific response can be enclosed and
independently checked before a production estimator is judged against it.
